% Generated draft: requires mathematical and source review.
\documentclass[11pt]{article}
\usepackage[margin=25mm]{geometry}
\usepackage{fontspec}
\setmainfont{Latin Modern Roman}
\usepackage{amsmath,amssymb,mathtools,unicode-math}
\setmathfont{Latin Modern Math}
\usepackage{enumitem,needspace,xurl,hyperref}
\setlength{\parindent}{0pt}
\setlength{\parskip}{5pt}
\setlength{\emergencystretch}{3em}
\newcommand{\N}{\mathbb{N}}
\newcommand{\Z}{\mathbb{Z}}
\newcommand{\Q}{\mathbb{Q}}
\newcommand{\R}{\mathbb{R}}
\newcommand{\C}{\mathbb{C}}
\newcommand{\F}{\mathbb{F}}
\newcommand{\A}{\mathbb{A}}
\newcommand{\PP}{\mathbb{P}}
\newcommand{\E}{\mathbb{E}}
\newcommand{\eps}{\varepsilon}
\newcommand{\dd}{\,\mathrm{d}}
\begin{document}
\section*{UnsolvedMath ID: 700755}
\textbf{Source problem code:} catalogue entry
\subsection{Definitions and mathematical statement}
A group is called a one-relator group if it admits a presentation \(G=\langle X\mid r\rangle\) with one defining relator. For a group \(G\), let
\[
T(G)=\{g\in G: g^m=1\text{ for some integer }m\geq 1\},
\qquad N=\langle T(G)\rangle.
\]
Thus \(N\) is the subgroup generated by all elements of finite order. It is normal, since conjugation preserves finite order. A group is free if it is isomorphic to the free group on some set; this includes the trivial group.

\textbf{Question (Kourovka Notebook 11.63).} For every one-relator group \(G\) with \(T(G)\neq\{1\}\), is it true that
\[
\text{for every subgroup }H\leq G,\qquad H\cap N=\{1\}\ \Longrightarrow\ H\text{ is free}?
\]
No finite generation assumption on \(H\) is imposed in the supplied statement.

\textbf{Supplied qualification.} The source record states that the answer is affirmative when \(G/N\) has non-trivial centre or satisfies a non-trivial identity. Here non-trivial centre means \(Z(G/N)\neq\{1\}\). A non-trivial identity means that there exist an integer \(k\geq1\) and a nonidentity word \(w\) in the free group on \(k\) generators such that
\[
w(q_1,\ldots,q_k)=1\qquad\text{for all }(q_1,\ldots,q_k)\in(G/N)^k.
\]
This affirmative special-case assertion is supplied source information, not an independently verified result. The record does not explicitly specify whether the generating set in its terminology ``one-relator group'' must be finite; the presentation definition above does not impose that additional restriction.
\subsection{Short English statement}
Suppose a one-relator group \(G\) contains a nonidentity element of finite order, and let \(N\) be generated by all its finite-order elements. Must every subgroup meeting \(N\) only in the identity be free? The supplied record labels the problem open and reports affirmative answers when \(G/N\) has non-trivial centre or satisfies a non-trivial identity; these assessments have not been independently verified.
\subsection{Sources}
\begin{itemize}
\item Supplied local UnsolvedMath registry record: permanent ID \(700755\), exact source code \texttt{catalogue entry}, title ``Kourovka Notebook Problem 11.63.'' Upstream revision: \texttt{d371ee917ac9647fbb06cd6fc529deec6f3dda82}. The record supplies the status ``open.''
\item \emph{The Kourovka Notebook (Unsolved Problems in Group Theory)}, original problem 11.63, 11th issue (1990), as identified in the supplied record with the 21st issue edition (2026), arXiv:1401.0300v48. The collection description names E. I. Khukhro and V. D. Mazurov as editors. This reference and the special-case assertion are supplied attribution and have not been independently checked.
\item Proposed in 1990; the supplied attribution begins ``D. I. Moldavanski'' and ends with a truncated TeX diacritic sequence. The exact supplied attribution is \texttt{D. I. Moldavanski\textbackslash{}u\{\textbackslash{}i}; its intended spelling remains unresolved.
\end{itemize}
\end{document}
